\subsection*{Learning a complete reading action}
For every source glyph we took its frozen question-aligned state, projected it to 64 dimensions and applied three residual convolutions of width five and dilations 1, 4 and 12. For a legal interval beginning at $i$ and ending at $j$, its score combined a learned start score, end score, length score conditioned on the start, and an eight-dimensional dot product between the two boundaries. We normalized over all contiguous intervals of one to 64 glyphs in the visible passage. When several intervals realized accepted responses, we minimized the negative logarithm of their summed probability. A separate binary cross-entropy trained the decision between the reading action and the free writer. At inference a positive decision selected the single highest-scoring interval and emitted its original pixels and stop. Otherwise the unchanged recurrent writer decoded freely.

We formed 5,187 training pairs with distinct pictured questions and disjoint acceptable answer images. For each pair we supplied the shortest sentence containing an annotated answer to each question, breaking ties in reading order. Both arms included ordinary DRCD question training, alternating the full passage with a reference-containing sentence across epochs. Each epoch included every shared-sentence training pair plus three additional presentations of the 190 person-question pairs. The control exchanged the two questions with probability one half in the paired presentations only. The primary evaluation retained all 864 pairs across 504 development passages, including 27 person pairs. Annotation-free evaluation selected a sentence by the previously fixed glyph-overlap rule. We separately compared 24 investigator-written role exchanges in simplified and traditional Chinese with two questions and both name orientations. We scored the entire image and stop, and for the text reference scored the complete name and stop from Qwen3-8B at temperature zero. The text reference did not train the visual model. An independent recount reconstructed the primary pairs from the original DRCD passages and accepted answer strings. The source-support reading was model-based, conducted before inspecting the fitted outputs on those pairs. It supplements the original exact scores.

For the action repair, DRCD, OCNLI, DuReader, training dictionary entries, and constructed forward and inverse questions provided mixed answers. A target representable by one contiguous source run trained the reading decision, and a target requiring independent writing trained the other decision. We updated only the decision network for 640 steps on 25,853 readable and 18,307 other training presentations. Every interval-scoring tensor was verified unchanged. The protected comparison includes descriptions whose required word order differs from the source, two bronze participant images, and human questions evaluated with the same inference-time evidence selection.
